\section{Conclusion}
\label{sec:conclusion}

A security monitoring and threat prevention platform was designed and implemented
for a Smart City Environmental Monitoring system built as a Service-Oriented
Architecture. Four threat classes are addressed: single-source HTTP flooding,
distributed botnet DDoS, a two-stage Advanced Persistent Threat, and two
independent email spam variants. Detection is achieved through thirteen custom
Wazuh rules combining frequency-based and cross-source correlation primitives.
Response is automated via a custom threat-blocker service that correctly handles
Docker NAT rewriting. Prometheus and Grafana provide real-time observability
across all threat classes, with pre-seeded counters eliminating cold-start blind
spots.

\subsection{Future Work}

Several directions would strengthen the platform in a production context.
The APT chain should be extended beyond credential capture to demonstrate
persistence establishment, data exfiltration, and lateral movement to downstream
services. The mail relay should be restricted to accept alert emails exclusively
from the alert service's internal IP, closing the trust gap identified in the
Evaluation. Inter-service communication should be upgraded from plain HTTP to
TLS with mutual certificate authentication. Finally, detection thresholds should
be calibrated against a baseline of real operator traffic before active response
is enabled, using a log-only monitoring period to quantify false positive rates
under legitimate load.

\subsection*{Resources}

The source code and configuration files are available at:
\url{https://github.com/odraude23/SSLE-SOA-Project1}.
